\hypertarget{chapter-16-c14-core-language-upgrades}{%
\section{CHAPTER 16: C14 CORE LANGUAGE
UPGRADES}\label{chapter-16-c14-core-language-upgrades}}

\hypertarget{c14-core-language-upgrades}{%
\section{C++14 CORE LANGUAGE
UPGRADES}\label{c14-core-language-upgrades}}

While C++11 was a revolution, C++14 was the ``refinement''
releasepolishing the rough edges of modern C++. It turned
\passthrough{\lstinline!constexpr!} from a toy into a powerful
compile-time engine and added ``quality of life'' features that brought
C++ syntax into the 21st century.

\hypertarget{relaxed-constexpr-the-deep-dive}{%
\subsection{1. Relaxed constexpr: The Deep
Dive}\label{relaxed-constexpr-the-deep-dive}}

In C++11, \passthrough{\lstinline!constexpr!} functions were strictly
functional: a single \passthrough{\lstinline!return!} statement, no
loops, no local variables. C++14 lifted these ``training wheels,''
allowing imperative logic.

\hypertarget{the-c11-vs.-c14-paradigm-shift}{%
\subsubsection{1.1 The C++11 vs.~C++14 Paradigm
Shift}\label{the-c11-vs.-c14-paradigm-shift}}

In C++11, you had to use recursion for almost everything. In C++14, you
can use standard algorithmic patterns.

\begin{lstlisting}[language={C++}]
// C++11: Functional/Recursive (Hard to read, heavy on stack during compilation)
constexpr int fib11(int n) {
    return (n <= 1) ? n : fib11(n - 1) + fib11(n - 2);
}

// C++14: Imperative/Iterative (Readable, efficient, familiar)
constexpr int fib14(int n) {
    if (n <= 1) return n;
    int a = 0, b = 1;
    for (int i = 2; i <= n; ++i) {
        int temp = a + b;
        a = b;
        b = temp;
    }
    return b;
}
\end{lstlisting}

\hypertarget{what-is-now-allowed}{%
\subsubsection{1.2 What is Now Allowed?}\label{what-is-now-allowed}}

\begin{itemize}
\tightlist
\item
  \textbf{Variable Declarations:} You can declare local variables
  (except \passthrough{\lstinline!static!} or
  \passthrough{\lstinline!thread\_local!}).
\item
  \textbf{Branching \& Loops:} \passthrough{\lstinline!if!},
  \passthrough{\lstinline!switch!}, \passthrough{\lstinline!for!},
  \passthrough{\lstinline!while!}, and
  \passthrough{\lstinline!do-while!} are all permitted.
\item
  \textbf{Mutation:} You can modify local variables within the function.
\item
  \textbf{Multiple Returns:} No longer restricted to a single
  expression.
\end{itemize}

\hypertarget{professional-note-the-constexpr-constraint}{%
\subsubsection{\texorpdfstring{1.3 Professional Note: The
\texttt{constexpr}
Constraint}{1.3 Professional Note: The constexpr Constraint}}\label{professional-note-the-constexpr-constraint}}

Even in C++14, a \passthrough{\lstinline!constexpr!} function cannot: 1.
Call non-\passthrough{\lstinline!constexpr!} functions. 2. Allocate
memory (until C++20). 3. Throw exceptions (though they can exist in
branches that are never taken at compile-time). 4. Use
\passthrough{\lstinline!asm!} blocks or \passthrough{\lstinline!goto!}
statements.

\begin{quote}
\textbf{Godhood Tip:} Use \passthrough{\lstinline!constexpr!} for any
computation that \emph{can} be done at compile-time. It doesn't just
save runtime; it allows the compiler to perform deeper optimizations on
the resulting constants.
\end{quote}

\hypertarget{binary-literals-digit-separators}{%
\subsection{2. Binary Literals \& Digit
Separators}\label{binary-literals-digit-separators}}

C++14 finally caught up with other languages by providing native binary
support and a way to make large numbers readable.

\hypertarget{hardware-level-clarity}{%
\subsubsection{2.1 Hardware-Level
Clarity}\label{hardware-level-clarity}}

For systems engineers and embedded developers, binary literals are a
godsend for bitmasking.

\begin{lstlisting}[language={C++}]
// Without Binary Literals (Hex/Octal required mental mapping)
uint8_t mask_old = 0x2A;

// With C++14 Binary Literals (Directly maps to hardware registers)
uint8_t mask_new = 0b0010'1010;

// Digit Separators (The single quote ')
// Can be placed anywhere to improve legibility
constexpr long double PLANCK_CONSTANT = 6.626'070'15e-34;
constexpr uint64_t MAX_CACHE_SIZE    = 0xFF'FF'FF'FF'FF'FF'FF'FF;
\end{lstlisting}

\hypertarget{the-deprecated-attribute}{%
\subsection{\texorpdfstring{3. The \texttt{{[}{[}deprecated{]}{]}}
Attribute}{3. The {[}{[}deprecated{]}{]} Attribute}}\label{the-deprecated-attribute}}

Standardizing how we tell other developers to stop using our old, broken
code.

\hypertarget{usage-patterns}{%
\subsubsection{3.1 Usage Patterns}\label{usage-patterns}}

The attribute can be applied to functions, classes, typedefs, variables,
and even namespaces.

\begin{lstlisting}[language={C++}]
namespace [[deprecated("Namespace is messy, use v2")]] LegacyAPI {
    struct [[deprecated]] OldData {
        int x;
    };
}

class Database {
public:
    [[deprecated("Use execute(Query&&) for better performance")]]
    void runRawSQL(const char* sql);
};
\end{lstlisting}

\textbf{Deep Dive:} Unlike \passthrough{\lstinline!\#pragma message!},
\passthrough{\lstinline![[deprecated]]!} is part of the language
standard. Compilers will emit a warning during the semantic analysis
phase, ensuring that the message is seen exactly when the deprecated
entity is utilized.
